\chapter{Multiplication III: Fermat FFT}
\label{ch:fermat}

A huge product can exceed the limits of a single wide CRT residue product.
One option is to cut each input into blocks, multiply every pair, and add the
overlapping results. BigMath also has an experimental path that transforms
the blocks first. It does fewer large inner products as the number of blocks
grows, at the cost of moving and mixing arrays of large values. The transform
is called \emph{outer} because the values it multiplies are themselves large
integers. BigMath uses its ordinary multiplier inside that step.

\begin{figure}[H]
\centering
\begin{tikzpicture}[node distance=5mm]
  \node[box,fill=sOrange!15] (in) {split into pieces\\and pad with zeros};
  \node[box,right=of in] (forward) {forward FFT\\of each input};
  \node[box,fill=sViolet!16,right=of forward] (multiply) {multiply pairs\\using BigMath};
  \node[box,right=of multiply] (inverse) {inverse FFT\\recovers sums};
  \node[box,fill=sOrange!15,right=of inverse] (out) {shift and add\\into result};
  \draw[arr] (in) -- (forward);
  \draw[arr] (forward) -- (multiply);
  \draw[arr] (multiply) -- (inverse);
  \draw[arr] (inverse) -- (out);
\end{tikzpicture}
\caption{The outer FFT rearranges the chunk products. Each inner product
is reduced modulo $2^n+1$ before the inverse transform.}
\end{figure}

\section{Why this modulus makes shifts cheap}

All the arithmetic inside the outer FFT uses remainders modulo $M=2^n+1$.
In this arithmetic, $2^n$ is one less than $M$, so it acts like $-1$.
If you shift a value and some bits spill over position $n$, split at
that position and subtract the high part from the low part:
\[
\text{low}+2^n\cdot\text{high}
   \equiv \text{low}-\text{high}\pmod{2^n+1}.
\]

\begin{figure}[H]
\centering
\begin{tikzpicture}[node distance=5mm]
  \node[box] (start) {$25\ll 3$\\$200$};
  \node[box,fill=sOrange!15,right=of start,minimum width=27mm] (high)
    {high part\\$3$};
  \node[box,fill=sAqua!13,right=1mm of high,minimum width=27mm] (low)
    {low 6 bits\\$8$};
  \node[box,right=8mm of low] (fold) {subtract high\\$8-3=5$};
  \draw[arr] (start) -- (high);
  \draw[arr] (low) -- (fold);
  \node[lbl,below=1mm of $(high.south)!0.5!(low.south)$] {$200=3\cdot64+8$};
\end{tikzpicture}
\caption{A shift by three bits followed by a fold, here modulo $65=2^6+1$.
If the subtraction is negative, add the modulus to get a nonnegative remainder.}
\end{figure}

Twiddle factors are the fixed multipliers in the butterfly steps
(\cref{ch:ntt}). The FFT chooses them to be powers of two, so applying one
is a shift and fold. A general pointwise product still requires the
ordinary multiplier. The implementation takes that full product and
folds its high half using the same subtraction. This style of exact FFT
uses integer remainders throughout. For the broader technique, see the
\href{https://gmplib.org/manual/FFT-Multiplication}{GNU MP manual on Fermat-style FFTs}.

\newpage
\section{A four-point product you can check by hand}

Multiply $13$ by $15$ in base 4, with one base-4 digit per piece. The inputs
are $(1,3)$ and $(3,3)$, low piece first. Choose four slots and fill the
unused input slots with zeros. The three sums that contribute to the
product occupy positions 0, 1, and 2, with room for all of them before
the transform's index would wrap back to 0.

\begin{figure}[H]
\centering
\begin{tikzpicture}[x=1.28cm,y=0.59cm]
  \foreach \i/\a/\b in {0/1/3,1/3/3,2/0/0,3/0/0} {
    \node[lbl] at (\i,1.55) {slot \i};
    \ifnum\i<2
      \node[box,minimum width=11mm,fill=sOrange!16] at (\i,0.75) {\a};
      \node[box,minimum width=11mm,fill=sAqua!16] at (\i,-0.25) {\b};
    \else
      \node[box,minimum width=11mm,fill=paper] at (\i,0.75) {\a};
      \node[box,minimum width=11mm,fill=paper] at (\i,-0.25) {\b};
    \fi
  }
  \node[lbl,anchor=east] at (-0.57,0.75) {$13$};
  \node[lbl,anchor=east] at (-0.57,-0.25) {$15$};
  \node[lbl,anchor=west] at (3.6,0.25) {zero padding leaves room for the product};
\end{tikzpicture}
\caption{The example uses four FFT slots. Their indices hold the powers of the piece
base after the inverse FFT recovers the coefficients.}
\end{figure}

Work modulo $65$. The four transform points are the powers of $8$:
$1,8,-1,-8$. They repeat after four steps because $8^2=64\equiv-1$ and
$8^4\equiv1$. Evaluate the two input recipes $1+3x$ and $3+3x$ at each point
and multiply the remainders.

\begin{center}
\small
\begin{tabular}{lrrrr}
\toprule
point & $1$ & $8$ & $-1$ & $-8$ \\
\midrule
$1+3x$ modulo $65$ & 4 & 25 & 63 & 42 \\
$3+3x$ modulo $65$ & 6 & 27 & 0 & 44 \\
pointwise product modulo $65$ & 24 & 25 & 0 & 28 \\
\bottomrule
\end{tabular}
\end{center}

At point 8, $25\cdot27=675=10\cdot64+35$. Folding gives $35-10=25$.
To undo the transform, weight each product by the inverse power of its
point, then divide the sum by four in the same modular arithmetic.
Here $4\cdot49\equiv1$, so division by four means multiplication by 49.
For example, the weights for coefficient 1 are $1,-8,-1,8$, giving
$24-8\cdot25+8\cdot28=48$ before scaling. These weights make four copies
of the coefficient we want. Contributions from the other coefficients
cancel, just as $1+8-1-8=0$.

\begin{center}
\small
\begin{tabular}{lrrrr}
\toprule
coefficient index & 0 & 1 & 2 & 3 \\
\midrule
weighted sum modulo $65$ & 12 & 48 & 36 & 0 \\
after multiplying by 49 modulo $65$ & 3 & 12 & 9 & 0 \\
\bottomrule
\end{tabular}
\end{center}

These are the integer coefficients. Each sums at most two products of
digits at most 3, so all are at most $18<65$. Now carry in base 4.
The $12$ at position 1 writes $0$ and sends $3$ to the $9$ at position 2.
That $12$ again writes $0$ and sends $3$ to position 3.
The resulting digits $(3,0,0,3)$ give $3+3\cdot4^3=195$.

\begin{figure}[H]
\centering
\begin{tikzpicture}[x=1.43cm,y=0.76cm]
  \foreach \i/\raw/\digit in {0/3/3,1/12/0,2/9/0,3/0/3} {
    \node[box,fill=sOrange!15,minimum width=11mm] (raw\i) at (\i,0) {\raw};
    \node[box,fill=sAqua!15,minimum width=11mm] (digit\i) at (\i,-1) {\digit};
    \draw[arr] (raw\i) -- (digit\i);
  }
  \node[lbl,anchor=east] at (-0.54,0) {coefficient};
  \node[lbl,anchor=east] at (-0.54,-1) {base-4 digit};
  \draw[arr,sOrange] (raw1.north) to[bend left=38] node[lbl,above] {$+3$} (raw2.north);
  \draw[arr,sOrange] (raw2.north) to[bend left=38] node[lbl,above] {$+3$} (raw3.north);
\end{tikzpicture}
\caption{Carry moves from low to high. A coefficient plus its incoming
carry supplies the digit below it and any carry to the next position.}
\end{figure}

\newpage
\section{How BigMath sets up a real product}
\label{sec:fermatparams}

Pieces in the library are whole runs of 64-bit limbs. Two choices make
the result exact.

\begin{enumerate}
  \item \textbf{Leave room for every position.} Choose a power-of-two number
        of slots $K=2^k$. With $a$ and $b$ limbs in the two operands, each
        piece holds $m=\ceil{(a+b)/(K-1)}$ limbs. The total number of output
        coefficients is one less than the sum of the input piece counts,
        and fits in $K$ slots.
  \item \textbf{Leave room for the value in each position.} A product of two
        $m$-limb pieces fits in $128m$ bits. Adding up to $K$ terms needs at
        most $k$ more bits. BigMath starts with $128m+k+1$ bits and rounds up
        to a multiple of $\max(64,K/2)$ for the modulus exponent $n$.
        That aligns residues to limb boundaries and makes $2^{2n/K}$ an
        available root for the FFT.
\end{enumerate}

The automatic selector starts at $k=6$ (64 slots) and takes the first fit
through $k=14$ with $n\le2^{28}$ bits. The forward and inverse FFTs share
the same shuffled slot order, saving a reordering pass. Large butterfly
layers use the pool. The caller does the pointwise products one at a time
so each inner multiplication can use that same pool. Each product reads the
significant limbs of its two residues in place. The caller can reuse the product
buffer at the next point after folding the result modulo $2^n+1$. For a square,
one forward transform serves both inputs. After the inverse FFT, add each
coefficient at its starting limb offset $im$, where $i$ is its position.
The coefficients can span several pieces, so the shifted additions
carry across their overlaps.

\section{When the dispatcher can use it}
\label{sec:fermatdispatch}

The build must use 64-bit limbs and the fused wide-CRT path. After both
CRT residue sets fail their size checks, the dispatcher
can try Fermat before splitting into overlap-add blocks
(\cref{sec:hugeproductdispatch}). It also checks the threshold
\code{BIGMATH\_NTT\_FERMAT\_MIN\_COEFFS}, counted in 32-bit product
coefficients, and checks whether the parameter selector found a fit.
The default threshold disables the path. See \file{FermatOuterMultiplication.h}.

\begin{table}[H]
\centering
\small
\begin{tabular}{lrrrr}
\toprule
Digits per input & \multicolumn{2}{c}{Warm time (s)} & \multicolumn{2}{c}{Peak RSS (GiB)} \\
\cmidrule(lr){2-3}\cmidrule(lr){4-5}
& Forced blocks & Fermat & Forced blocks & Fermat \\
\midrule
250 million & 3.39 & 2.88 & 1.49 & 1.71 \\
2.5 billion & 27.66 & 29.12 & 10.26 & 14.34 \\
5 billion & 55.65 & 77.79 & 19.30 & 20.12 \\
\bottomrule
\end{tabular}
\caption{Local comparisons on a 32\,GiB Apple M1 Max with up to eight
threads. Both modes were forced at these sizes. The forced-block control used
three blocks per input. Warm products matched GMP hashes. The 5-billion-digit
Fermat run incurred system swap.}
\end{table}

The crossover beyond the wide-CRT limits still needs measurements at
those sizes. The default build handles the balanced 2.5- and 5-billion-digit
cases in the table with single wide residue products. Full measurements and
reproduction commands are in \file{BENCHMARK.md}.
